\documentclass[11pt]{article}
\usepackage[utf8]{inputenc}
\usepackage[T1]{fontenc}
\usepackage{lmodern}
\usepackage[margin=1in]{geometry}
\usepackage{amsmath,amssymb,mathtools}
\usepackage{graphicx}
\usepackage{grffile}
\usepackage{float}
\usepackage{hyperref}
\usepackage{parskip}
\usepackage{enumitem}
\usepackage{xurl}
\setlength{\parskip}{0.65em}
\setlength{\parindent}{0pt}
\begin{document}
\title{Daily Research Digest --- 2026-04-14}
\author{}
\date{}
\maketitle
\textbf{Topics:} galaxy clusters, galaxies, gravitational lensing, dark matter.
\section*{Synthesis}
\_Digest generation failed: timed out\_
\newpage
\section*{Paper Summaries and Figures}
\subsection*{1. Cosmological inference with halo clustering reconstructed from the redshift-space galaxy distribution}
\textbf{Focus:} galaxies\\
\textbf{Authors:} Ryuichiro Hada, Teppei Okumura\\
\textbf{Published:} 2026-04-13T16:33:23+00:00\\
\textbf{PDF:} \href{https://arxiv.org/pdf/2604.11694v1}{https://arxiv.org/pdf/2604.11694v1}\\
\textbf{Summary:}
Accurate modeling of small-scale redshift-space clustering is crucial for full shape RSD analyses, where satellite galaxies contribute to 1-halo terms and Finger-of-God distortions. We investigate halo reconstruction based on the cylinder grouping (CG) method of Okumura et al. (2017), which selects an effective halo center tracer from the observed galaxy distribution, and how it impacts cosmological parameter inference. Using DESI-like luminous red galaxy mock catalogs from the AbacusSummit simulations at $z=1.1$, we perform effective field theory (EFT)-based full-shape modeling of the power spectrum of the reconstructed-halo sample. We show that the dominant reconstruction-induced systematics can be described and incorporated within the standard EFT framework. In particular, a simple multipole-dependent rescaling inferred directly from the data on large scales captures the dominant effect, while residual small-scale changes are absorbed by the standard counterterm and stochastic sector, without introducing additional reconstruction-specific parameters. The reconstructed-halo sample yields unbiased constraints on cosmological parameters, including the growth rate $f\sigma _8$ and Alcock-Paczynski parameters. Compared to the galaxy sample, it enables both improved robustness and increased statistical precision: the inferred $f\sigma _8$ remains stable when extending the fit beyond $k_{\max}\simeq 0.2\,h\,{\rm Mpc}^{-1}$, with its uncertainty reduced by more than $20\%$.
\newpage
\subsection*{2. Galactic Diffuse Gamma-Ray and Neutrino Emission from Cosmic-Ray Interactions in Stellar Atmospheres}
\textbf{Focus:} dark matter\\
\textbf{Authors:} Yanbo Wang, Zhenglong Wang, Rui Zhang, Yi Zhang\\
\textbf{Published:} 2026-04-13T15:23:04+00:00\\
\textbf{PDF:} \href{https://arxiv.org/pdf/2604.11617v1}{https://arxiv.org/pdf/2604.11617v1}\\
\textbf{Summary:}
The Galactic diffuse gamma-ray emission is conventionally modeled as the product of cosmic-ray interactions with the interstellar medium. However, the cumulative contribution of stellar atmospheres acting as hadronic interaction targets remains an unexplored multi-messenger background. In this work, we present the first systematic evaluation of this stellar diffuse emission by coupling MESA stellar evolution profiles and magnetic-field-modulated cosmic-ray transport with a 3D Galactic population synthesis framework. We find that the cumulative stellar contribution to the Galactic diffuse gamma-ray flux is negligible at 1 TeV, and the associated diffuse neutrino flux ($\sim 10^{-16}\;\mathrm{TeV\;cm^{-2}\;s^{-1}\;sr^{-1}}$) remains orders of magnitude below current IceCube limits. Nevertheless, at ultra-high energies ($>10\;\mathrm{TeV}$), this emission establishes an irreducible local background that overtakes the strongly attenuated extragalactic isotropic gamma-ray background. Our results demonstrate that the Galactic stellar ensemble is a strictly sub-dominant background, indicating that stellar subtraction templates are not required for identifying Galactic PeVatrons or constraining dark matter annihilation.
\subsubsection*{Selected figures}
\begin{figure}[H]
\centering
\includegraphics[width=\linewidth,height=0.42\textheight,keepaspectratio]{/home/kf/research-assistant/data/cache/figures/http-arxiv.org-abs-2604.11617v1/figure\_1.png}
\end{figure}
\newpage
\subsection*{3. Extragalactic microlensing through Ultra Diffuse Galaxies}
\textbf{Focus:} galaxies, dark matter\\
\textbf{Authors:} Sung Kei Li, Thomas Broadhurst, Jose M. Diego, Jeremy Lim, Jose M. Palencia, James Nianias\\
\textbf{Published:} 2026-04-13T12:10:29+00:00\\
\textbf{PDF:} \href{https://arxiv.org/pdf/2604.11368v1}{https://arxiv.org/pdf/2604.11368v1}\\
\textbf{Summary:}
Stellar microlensing is a powerful method to constrain compact dark matter models, uncover binary stars, and exoplanets during caustic crossing events. At cosmological distances, \{\textbackslash{}it James-Webb Space Telescope\} (\{\textbackslash{}it JWST\}) is routinely detecting microlensed giant stars in highly magnified galaxies behind massive lensing clusters. Here, we explore for the first time microlensing in modest redshift galaxies commonly seen through local Ultra Diffuse Galaxies (UDGs). Using the UDG NGC1052-DF2 as a case study, we found that detecting UDG microlensing events through UDGs is possible. However, a low total UDG microlensing event rate of $\sim 5.6\times10^{-2}\,\textrm{yr}^{-1}$ over its five background galaxies is expected for typical \{\textbackslash{}it JWST\} $\sim 29\,$mag visits, and a low Vera Rubin Legacy Survey of Space and Time (LSST) detection rate of $\sim 2\times10^{-8}\,\textrm{yr}^{-1}$ such that NGC1052-DF2 might not be a prime target given its lack of low-redshift background galaxies. \{\textbackslash{}it Euclid\} is ideal for identifying samples of low-redshift star-forming galaxies seen through local galaxies for deeper cadenced follow-up, where our zeroth-order calculation estimates that $\mathcal{O}(1-10)$ events per year are expected over the whole sky under the monitoring of LSST. Finally, we postulate that UDG microlensing will allow an independent estimate of the initial mass function (IMF) and the stellar multiplicity in the low mass regime, of considerable interest for UDG galaxies, where stellar mass has been claimed to predominate over dark matter in some cases, including NGC1052-DF2.
\subsubsection*{Selected figures}
\begin{figure}[H]
\centering
\includegraphics[width=\linewidth,height=0.42\textheight,keepaspectratio]{/home/kf/research-assistant/data/cache/figures/http-arxiv.org-abs-2604.11368v1/figure\_1.png}
\end{figure}
\newpage
\subsection*{4. Charged Black Holes in KR-gravity Surrounded by Perfect Fluid Dark Matter}
\textbf{Focus:} dark matter\\
\textbf{Authors:} Faizuddin Ahmed, Mohsen Fathi, Edilberto O. Silva\\
\textbf{Published:} 2026-04-13T11:53:25+00:00\\
\textbf{PDF:} \href{https://arxiv.org/pdf/2604.11357v1}{https://arxiv.org/pdf/2604.11357v1}\\
\textbf{Summary:}
In this work, we systematically investigate the null geodesics of electrically charged black holes in a gravitational framework that incorporates Lorentz violation induced by a background Kalb-Ramond (KR) field, in the presence of perfect-fluid dark matter. The properties of the photon sphere, black hole shadow, and photon trajectories are analyzed in detail. Furthermore, to explore the combined effects of Lorentz violation and dark matter on the motion of neutral test particles, we examine the innermost stable circular orbit (ISCO) in this spacetime. In addition, the epicyclic frequencies of test particles are studied to gain further insight into the dynamical behavior of particle motion around these black holes. The main analytical results are complemented by a phenomenological QPO analysis, a thermodynamic investigation, and a discussion of the sparsity of Hawking radiation, allowing us to connect optical, dynamical, and thermodynamic signatures within a single framework.
\subsubsection*{Selected figures}
\begin{figure}[H]
\centering
\includegraphics[width=\linewidth,height=0.42\textheight,keepaspectratio]{/home/kf/research-assistant/data/cache/figures/http-arxiv.org-abs-2604.11357v1/figure\_1.png}
\end{figure}
\newpage
\subsection*{5. A consistent MOND modelling of the Bullet Cluster}
\textbf{Focus:} dark matter\\
\textbf{Authors:} X. Hernandez\\
\textbf{Published:} 2026-04-12T20:46:08+00:00\\
\textbf{PDF:} \href{https://arxiv.org/pdf/2604.10811v1}{https://arxiv.org/pdf/2604.10811v1}\\
\textbf{Summary:}
It is a common miss-conception that 1E 0657-56, the "Bullet Cluster", is somehow inconsistent with MOND expectations. The argument centres on the fact that the baryonic matter distribution of this system is dominated by the X-ray emitting gas, while the total projected surface density required under General Relativity to explain the observed lensing signal, centres on the observed galaxies. This is sometimes interpreted as being in conflict with MOND, as under such an interpretation, it is naively assumed that all dark matter being absent, the gravitational potential should necessarily be dominated by the largest mass distribution, that of the gas. However, just as under General Relativity, under MOND, the total gravitational potential of a system depends sensitively upon the volume density and not just on the total mass. It is shown in this \{\textbackslash{}it letter\} that the surface density which QUMOND predicts will be inferred under a standard gravity framework from the total gravitational potential of the Bullet Cluster, closely matches what General Relativity inferences of lensing observations return. The close-to-point-like galaxies imply under QUMOND a relatively much larger surface density signal than what is expected from the Mpc scale gas distribution.
\newpage
\subsection*{6. Reionization Topology as a Probe of Self-Interacting Dark Matter}
\textbf{Focus:} dark matter\\
\textbf{Authors:} Zihan Wang\\
\textbf{Published:} 2026-04-12T16:59:09+00:00\\
\textbf{PDF:} \href{https://arxiv.org/pdf/2604.10726v1}{https://arxiv.org/pdf/2604.10726v1}\\
\textbf{Summary:}
We introduce a framework connecting dark matter self-interactions (SIDM) to the large-scale topology of cosmic reionization in this paper. SIDM core formation reduces the gas binding energy in high-$z$ halos, enhancing supernova-driven clearing of ionizing-photon escape channels. We decompose the observable signatures into two scale-dependent levers a percent-level shift in the emissivity-weighted halo bias $b_\gamma $ that modifies large-scale 21\textbackslash{},cm power, and a factor of 2--4 suppression of emissivity shot noise from increased duty cycles that reshapes intermediate-scale ionization morphology. We derive analytic predictions for the 21\textbackslash{},cm power spectrum ratio and validate them with a halo-by-halo semi-numerical excursion-set framework at $128^3$ resolution where individual halo duty-cycle stochasticity is resolved. For $\sigma /m = 1$--$10\;\mathrm{cm^2/g}$, we find $\sim 60$--$80\%$ suppression of emissivity shot-noise power, a $12$--$21\%$ reduction in the emissivity variance ratio $\langle \dot{N}^2\rangle/\langle\dot{N}\rangle^2$, and a $50$--$110\%$ increase in the Euler characteristic of the ionization field at fixed $\bar{x}_{\rm HI} = 0.5$. SIDM produces more numerous, more uniformly distributed HII bubbles compared to CDM's topology of fewer large bubbles around rare bright sources. Theintermediate scale signatures are potentially detectable by SKA1-Low in $\sim 1000$ hours. Our results establish reionization topology as a new, complementary probe of dark matter microphysics at mass scales $M \sim 10^{10}$--$10^{11}\,M_\odot$ and redshifts $z \sim 6$--$10$.
\subsubsection*{Selected figures}
\begin{figure}[H]
\centering
\includegraphics[width=\linewidth,height=0.42\textheight,keepaspectratio]{/home/kf/research-assistant/data/cache/figures/http-arxiv.org-abs-2604.10726v1/figure\_1.png}
\end{figure}
\newpage
\subsection*{7. Strong gravitational lensing and Quasiperiodic oscillations as a probe for an electrically charged Lorentz symmetry-violating black hole}
\textbf{Focus:} gravitational lensing\\
\textbf{Authors:} Sohan Kumar Jha\\
\textbf{Published:} 2026-04-12T10:35:32+00:00\\
\textbf{PDF:} \href{https://arxiv.org/pdf/2604.10572v1}{https://arxiv.org/pdf/2604.10572v1}\\
\textbf{Summary:}
This study examines the combined effect of electric charge and Lorentz symmetry breaking (LSB) on the observables of strong gravitational lensing (SGL) and the dynamics of quasiperiodic oscillations (QPOs) around an electrically charged, Lorentz symmetry-violating (LV) black hole (QKR BH). We first explore the SGL, which unravels an interesting effect that the two combined generate. We find cases where the competing effect of charge and LV cancels each other, leaving the underlying quantity unchanged from that of a \textbackslash{}s BH. We find bounds on the LV parameter utilizing observations related to the shadow angular size of supermassive black holes (SMBHs) $M87^*$ and $SgrA^*$. No bound could be gleaned for the charge from these shadow observations. Observations of QPOs in microquasars provide an alternative method to probe our model and to extract bounds on its parameters. We use experimental data for the microquasars $GRO J1655-40$ and $XTE J1550-564$. Here we obtain bounds on both parameters. Our results provide deeper insights into the interplay between charge and LSB in the strong-gravity regime.
\subsubsection*{Selected figures}
\begin{figure}[H]
\centering
\includegraphics[width=\linewidth,height=0.42\textheight,keepaspectratio]{/home/kf/research-assistant/data/cache/figures/http-arxiv.org-abs-2604.10572v1/figure\_1.png}
\end{figure}
\newpage
\subsection*{8. Search for dark sector and rare decays at BESIII}
\textbf{Focus:} dark matter\\
\textbf{Authors:} Zhi-Jun Li, Zheng-Yun You\\
\textbf{Published:} 2026-04-12T01:50:29+00:00\\
\textbf{PDF:} \href{https://arxiv.org/pdf/2604.10405v1}{https://arxiv.org/pdf/2604.10405v1}\\
\textbf{Summary:}
The BESIII experiment has collected a large data sample of charmonium, charm mesons, hyperons, and other light mesons. These data provide a unique opportunity to explore the dark sector and rare decays. We present recent dark sector results from the BESIII experiment, including searches for sub-GeV dark matter in $η\to \pi ^0 + \text{invisible}$ and $J/ψ\to \phi + \text{invisible}$, searches for dark baryon particles in $Ξ^- \to \pi ^- +\rm{invisible}$, and searches for light vector bosons in $χ_{cJ} \to J/ψX$, where $X \to e^+ e^-$. In addition, we also present the recent rare decay searches, including a search for $ψ\to e \mu $, as well as searches for several baryon number and lepton number violation processes, and searches for various charmonium weak decays.
\newpage
\subsection*{9. DREAMuS: Dark matter REsearch with Advanced Muon Source}
\textbf{Focus:} dark matter\\
\textbf{Authors:} Xiang Chen, Zejia Lu, Liangwen Chen, Jun Gao, Shao-Feng Ge, Zhanxu Hao, Yang Hu, Bingzhi Li, Cen Mo, Zhiyu Sun, Huayang Wang, Chonghao Wu, Yu Xu, Xueheng Zhang, Yulei Zhang, Liang Li\\
\textbf{Published:} 2026-04-11T15:49:07+00:00\\
\textbf{PDF:} \href{https://arxiv.org/pdf/2604.10257v1}{https://arxiv.org/pdf/2604.10257v1}\\
\textbf{Summary:}
We propose DREAMuS, a fixed-target experiment at the High Intensity Heavy-Ion Accelerator Facility (HIAF), to search for muon-philic dark matter mediated by light flavor-violating bosons. DREAMuS is designed to probe the parameter space of a muon-philic dark matter (DM) mediated by a light flavor-violating boson, specifically a vector $Z'$ (or a scalar $\phi $) which is produced in muon-nucleus interactions and decays into dark matter particles with a distinctive detector signature. Precision tracking and time-of-flight measurements are used to suppress the Standard Model backgrounds. We find that DREAMuS can achieve competitive sensitivity in the GeV-scale muon-philic dark matter parameter space, reaching sensitivity to couplings at the $10^{-4}$, especially in the few-hundred-MeV region.In addition to a $\mu ^-$ run, we highlight the potential of a complementary $\mu ^+$ beam option, further improving sensitivity to dark matter below 200 $\mathrm{MeV}$ by an order of magnitude.
\newpage
\subsection*{10. Do little red dots really form a distinct class of astronomical objects?}
\textbf{Focus:} galaxies\\
\textbf{Authors:} Jean-Baptiste Billand, David Elbaz, Maximilien Franco, Fabrizio Gentile, Emanuele Daddi, Mauro Giavalisco, Dale D. Kocevski, Joseph S. W. Lewis, Benjamin Magnelli, Valentina Sangalli, Maxime Tarrasse\\
\textbf{Published:} 2026-04-13T16:22:12+00:00\\
\textbf{PDF:} \href{https://arxiv.org/pdf/2604.11677v1}{https://arxiv.org/pdf/2604.11677v1}\\
\textbf{Summary:}
JWST observations have identified a class of enigmatic sources known as "Little Red Dots" (LRDs). These have been interpreted as a distinct class of active galactic nuclei (AGN) and host galaxies, potentially involving "quasi-stars" or Black Hole stars (BH*). However, two questions remain: is there a clear discontinuity between LRDs and field galaxies, and do LRDs form a homogeneous population? In this work, we address these issues by introducing a continuous metric to evaluate the "LRDness" of galaxies. We measure their compactness ($\delta _{compact}$), the sharpness of the V-shaped spectral energy distribution ($\delta _{v-shape}$), and the strength of the broad Balmer line emission. We apply this approach to a sample of \textasciitilde{}48,000 galaxies with photometric and \textasciitilde{}5,000 with spectroscopic information, selected over \textasciitilde{}750 arcmin\textasciicircum{}2. We find that V-shape prominence correlates strongly with morphology without a clear transition at common LRD selection thresholds: the fraction of compact galaxies rises continuously with V-shape intensity. Similarly, broad H$\alpha $ strength increases with both V-shape sharpness and compactness. The [N II] deficit is not an exclusive feature of LRDs but a global property of compact, metal-poor galaxies. Only the 3\% most extreme LRDs present a prominent Balmer break (>3) of potentially non-stellar origin. LRDs and non-LRDs follow similar trends in the evolution of the Balmer decrement with V-shape sharpness, suggesting a shared physical origin, likely dust attenuation. Estimated dust masses (\textasciitilde{}4-7 x 10\textasciicircum{}4 M\_\{sun\}) and luminosities are low enough to account for their non-detection by ALMA. We conclude that most LRDs do not represent a separate class of objects, but rather the extreme tail of a continuous distribution of galaxies and broad H$\alpha $ emitters, consistent with a classical broad line region and dust attenuation.
\end{document}
